\documentclass[11pt]{amsart}

\usepackage[T1]{fontenc}
\usepackage{lmodern}
\usepackage{microtype}
\usepackage{amsmath,amssymb,amsthm}
\usepackage[hidelinks]{hyperref}

\hypersetup{
  pdftitle={Counterexamples to the Formula g(C3 plus C3p) = 3p + 3}
}

\newtheorem{theorem}{Theorem}
\newtheorem{corollary}[theorem]{Corollary}
\theoremstyle{remark}
\newtheorem{remark}[theorem]{Remark}

\newcommand{\gH}{\mathsf{g}}

\title[Counterexamples to a Harborth formula]
{Counterexamples to the Formula
\texorpdfstring{$\gH(C_3\oplus C_{3p})=3p+3$}
{g(C3 plus C3p) = 3p + 3}}
\date{}
\subjclass[2020]{11B30, 20K01}
\keywords{Harborth constant, zero-sum subset, finite abelian group}

\begin{document}

\begin{abstract}
For every positive integer $n\equiv5\pmod6$, we construct a
$(3n+3)$-element subset of $C_3\oplus C_{3n}$ containing no zero-sum
subset of cardinality $3n$. Hence
\[
  \gH(C_3\oplus C_{3n})\ge 3n+4.
\]
For every prime $p\equiv5\pmod6$, this contradicts Theorem~3.1 of
Guillot, Marchan, Ordaz, Schmid, and Zerdoum, which asserts
$\gH(C_3\oplus C_{3p})=3p+3$ for all primes $p\ne3$.
The smallest counterexample produced by this construction is $p=5$. We
also identify the erroneous step in the proof of their
Proposition~3.7.
\end{abstract}

\maketitle

Let $C_m$ denote the cyclic group of order $m$, written additively.
For a finite abelian group $G$, its Harborth constant $\gH(G)$ is the
least integer $k$ such that every subset of $G$ of cardinality at least
$k$ contains a zero-sum subset of cardinality $\exp(G)$; see
Bajnok~\cite[Chapter~F.3]{Bajnok}. Guillot et
al.~\cite[Theorem~3.1]{GuillotEtAl} state that
\[
  \gH(C_3\oplus C_{3p})=3p+3
\]
for every prime $p\ne3$. The following construction disproves this
formula for every prime $p\equiv5\pmod6$ and applies more generally to
all integers in that congruence class.

\begin{theorem}\label{thm:main}
If $n\equiv5\pmod6$, then
\[
  \gH(C_3\oplus C_{3n})\ge 3n+4.
\]
\end{theorem}

The bound $\gH(C_3\oplus C_{3n})\ge3n+3$ for $n\ge2$ is due to Kiefer;
see also \cite[Prop.~F.104]{Bajnok}. Theorem~\ref{thm:main} improves it
by one for $n\equiv5\pmod6$.

\begin{proof}
Write $n=3m-1$, so $m=(n+1)/3$ is even. Define
\[
  \rho:C_3\oplus C_{3n}\longrightarrow C_n,
  \qquad \rho(a,b)=b\pmod n,
\]
and, using the integer representatives in $\{0,\ldots,n-1\}$, set
\[
  J=\{m,m+1,\ldots,2m-1\}\subset C_n,
  \qquad A_n=\rho^{-1}(J).
\]
Every fiber of $\rho$ has cardinality $9$, and therefore
\[
  |A_n|=9m=3n+3.
\]
For $j\in J$ one has
\[
  \rho^{-1}(j)
  =\{(a,j+\ell n):a,\ell\in\{0,1,2\}\},
\]
where the second coordinate is taken in $C_{3n}$. Hence the sum of this
fiber is $(0,9j)$, and
\[
  \sigma(A_n)
  =\left(0,9\sum_{j=m}^{2m-1}j\right)
  =\left(0,\frac{9mn}{2}\right)
  =0
  \quad\text{in }C_3\oplus C_{3n},
\]
since $m$ is even and $9mn/2$ is divisible by $3n$.

Suppose that $B\subset A_n$ has cardinality $3n$ and sum $0$, and put
$R=A_n\setminus B$. Then $|R|=3$ and
\[
  \sigma(R)=\sigma(A_n)-\sigma(B)=0.
\]
For each $r\in R$, let $\widetilde{\rho(r)}\in J$ be the chosen integer
representative of $\rho(r)$. Then
\[
  n+1=3m
  \le \sum_{r\in R}\widetilde{\rho(r)}
  \le 3(2m-1)=2n-1.
\]
This interval contains no multiple of $n$, so
$\rho(\sigma(R))\ne0$, a contradiction. Thus $A_n$ contains no
zero-sum subset of cardinality $3n$. Since
$\exp(C_3\oplus C_{3n})=3n$, the claimed lower bound follows.
\end{proof}

\begin{corollary}\label{cor:consequences}
For every prime $p\equiv5\pmod6$,
\[
  \gH(C_3\oplus C_{3p})\ge3p+4.
\]
In particular,
\[
  \gH(C_3\oplus C_{15})\ge19,
\]
contradicting the value $18$ asserted in
\cite[Theorem~3.1]{GuillotEtAl}.
Moreover, taking $n=35$ gives
\[
  \gH(C_3\oplus C_{105})\ge109,
\]
which disproves the suggested extension of the formula to composite
parameters in \cite[p.~632]{GuillotEtAl}.
\end{corollary}

The construction establishes these lower bounds only; it does not determine
the exact Harborth constants in the affected cases.

\begin{remark}[Failure of Proposition~3.7]\label{rem:gap}
Let $p\equiv5\pmod6$ be prime and write $p=3m-1$. In the decomposition
$C_3\oplus C_{3p}=H_1\oplus H_2$ used in \cite{GuillotEtAl}, one has
$H_1=\ker\rho\cong C_3^2$, and the restriction of $\rho$ identifies
$H_2$ with $C_p$. Viewing the set $A_p$ above as a squarefree sequence,
every $H_1$-fiber has $H_2$-coordinate set
\[
  J=\{m,m+1,\ldots,2m-1\}.
\]
Moreover, its $H_1$-support is all of $H_1$ and
\[
  \sigma(\pi_1(A_p))=m\sum_{h\in H_1}h=0.
\]
Thus $A_p$ satisfies the hypotheses of
\cite[Proposition~3.7]{GuillotEtAl}, while Theorem~\ref{thm:main}
shows that it violates the conclusion.

The error occurs near the end of the proof of that proposition. For an
arithmetic progression
\[
  P_h=\{s_he,(s_h+1)e,\ldots,(s_h+m-1)e\}\subset C_p,
\]
write $\Sigma_3(P_h)$ for the set of sums of three distinct terms.
Its elements have integer coefficients in
\[
  [3s_h+3,\,3s_h+p-5].
\]
The proof infers from $0\notin\Sigma_3(P_h)$ that the upper endpoint is
less than $p$. This overlooks the possibility that the entire interval
lies strictly between $p$ and $2p$. For the fibers above, $s_h=m$, and
for $p\ge11$ the interval is
\[
  [p+4,\,2p-4],
\]
which contains no multiple of $p$. For $p=5$ one has $m=2$, so each
progression $P_h$ has only two terms and $\Sigma_3(P_h)$ is empty.
Proposition~3.7 is used in the proof of
\cite[Theorem~3.1]{GuillotEtAl}.
\end{remark}

\begin{thebibliography}{9}

\bibitem{Bajnok}
B.~Bajnok,
\emph{Additive Combinatorics: A Menu of Research Problems},
Discrete Mathematics and Its Applications,
CRC Press, Boca Raton, 2018.

\bibitem{GuillotEtAl}
P.~Guillot, L.~E.~Marchan, O.~Ordaz, W.~A.~Schmid, and H.~Zerdoum,
On the Harborth constant of $C_3\oplus C_{3p}$,
\emph{J. Th\'eor. Nombres Bordeaux}
\textbf{31} (2019), no.~3, 613--633,
\href{https://doi.org/10.5802/jtnb.1097}{doi:10.5802/jtnb.1097}.

\end{thebibliography}

\end{document}